\documentclass[11pt]{article}

\usepackage[utf8]{inputenc}
\usepackage[T1]{fontenc}
\usepackage[margin=1in]{geometry}
\usepackage{booktabs}
\usepackage{amsmath}
\usepackage{listings}
\usepackage{xcolor}
\usepackage{hyperref}
\hypersetup{colorlinks=true, linkcolor=blue!50!black, citecolor=blue!50!black,
            urlcolor=blue!50!black}

% GENERATED by paper/make_tables.py -- do not edit by hand.

\newcommand{\AcCandidates}{45}
\newcommand{\AcScreenedOut}{32}
\newcommand{\AcRepos}{13}
\newcommand{\AcTools}{158}
\newcommand{\AcCovered}{126}
\newcommand{\AcUncovered}{32}
\newcommand{\AcUncoveredPct}{20.3\%}
\newcommand{\AcRepoUntested}{1}
\newcommand{\AcRepoUntestedPct}{7.7\%}
\newcommand{\AcSensRepos}{7}
\newcommand{\AcSensTools}{87}
\newcommand{\AcSensUncovered}{18}
\newcommand{\AcSensPct}{20.7\%}
\newcommand{\AcRunDate}{2026-08-04}
\newcommand{\AcUntestedRepoList}{aiwaves-cn/agents}
\newcommand{\AcCiLo}{14.7}
\newcommand{\AcCiHi}{27.2}
\newcommand{\AcConcPct}{84.4\%}
\newcommand{\AcVcsPct}{46.7\%}
\newcommand{\AcShellPct}{33.3\%}
\newcommand{\AcShellN}{36}


% Long \texttt{} paths (part_b/demo_double_refund.py) cannot be hyphenated, so TeX
% would rather overrun the margin than stretch a line. Let it stretch instead.
\emergencystretch=3em
\hbadness=2000

\lstset{
  basicstyle=\ttfamily\footnotesize,
  breaklines=true,
  frame=single,
  framesep=4pt,
  rulecolor=\color{black!30},
  columns=fullflexible,
}

\title{\textbf{Untested by Construction: Test Coverage of Write-Capable Tools\\
in Open-Source LLM Agents}}

\author{Rithik Satarla\\
\texttt{rithiksatarla@gmail.com}}

\date{\today}

\begin{document}
\maketitle

\begin{abstract}
Large language model agents increasingly hold credentials for systems that change
state: they issue refunds, delete records, push branches, and send mail. The dominant
way to test such an agent is to record an HTTP response and replay it. A replay
fixture, however, has no state, and so it cannot express the failure mode specific to
write operations---an operation applied twice, or applied when it should have been
refused. An agent that retries a refund it already completed passes a fixture-based
suite and double-refunds in production.

We ask how much of the open-source agent ecosystem tests its write paths at all. We
screened \AcCandidates{} Python-primary agent repositories, mechanically extracted
framework-idiomatic tool definitions, and retained the \AcTools{} definitions whose
bodies perform a detectable state mutation across the \AcRepos{} repositories that had
any. Using a deliberately generous coverage criterion---a tool is covered if its
identifier appears anywhere in the repository's test sources, with no assertion or
execution required---we find that \AcUncoveredPct{} of write-capable tool definitions
(\AcUncovered{} of \AcTools{}) are never referenced by any test. Restricting to the
\AcSensRepos{} repositories whose entire test suite was read gives \AcSensPct{},
so the estimate does not depend on sampling. Because the criterion is an upper bound
on real coverage, \AcUncoveredPct{} is a lower bound on what goes untested.

We then describe AgentCage, a stateful alternative: a passive HTTP interceptor plus an
in-memory API model in which writes mutate state and subsequent reads observe the
mutation. We pre-register an evaluation of whether this surfaces defects that replay
fixtures cannot. All code, data, and the detector revision history are public.
\end{abstract}

\section{Introduction}

An LLM agent with tool access is, operationally, a program that issues requests against
external systems. A growing fraction of those requests are writes. Agent frameworks ship
tools that commit to git, delete cloud objects, run shell commands, execute SQL, send
email, and move money.

The testing practice these agents inherited comes from the world of read-mostly API
clients: record a response, replay it. This works for a tool that fetches a document.
It fails for a tool that moves money, and it fails in a specific, structural way.
Consider an agent that issues a refund, fails to parse the response, and retries:

\begin{lstlisting}
refund(charge_id, 5000)   # fixture replies 200 OK
refund(charge_id, 5000)   # fixture replies 200 OK  <- suite passes
                          # Stripe would reply 400 charge_already_refunded
\end{lstlisting}

The fixture answers identically however many times it is asked, because it holds no
state. The defect lives entirely in the state transition, so no assertion over the
fixture's responses can detect it. The test suite is not merely incomplete; it is
incapable of expressing the property.

This paper makes two contributions.

\textbf{A measurement (Section~\ref{sec:parta}).} We quantify how much of the
open-source agent ecosystem tests its write-capable tools, using a mechanical,
reproducible procedure over \AcCandidates{} repositories. The headline is
\AcUncoveredPct{} of \AcTools{} write-capable tool definitions never referenced by any
test. We report the result together with the seven revisions our detectors underwent,
two of which removed false positives that would have made the finding look worse.

\textbf{An artefact and a pre-registered evaluation (Section~\ref{sec:partb}).} We
describe AgentCage: passive HTTP capture plus a stateful API model. We state in advance
what would count as evidence that it surfaces defects fixtures cannot, and what would
falsify that claim. The evaluation has not been run; no results are reported for it.

\paragraph{Scope.} This is a study of \emph{test coverage}, not of agent safety in
deployment. We do not claim the uncovered tools are buggy, nor that any agent has
caused harm. We claim only that a fifth of the write-capable tools we could identify
are not touched by their projects' tests, under the weakest definition of ``touched''
we could construct.

\section{Related work}

Agent benchmarks evaluate whether agents \emph{accomplish} tasks. ReAct~\cite{react}
established interleaved reasoning and acting; Toolformer~\cite{toolformer},
Gorilla~\cite{gorilla}, and ToolLLM~\cite{toolllm} study tool invocation.
SWE-bench~\cite{swebench}, WebArena~\cite{webarena}, AgentBench~\cite{agentbench}, and
GAIA~\cite{gaia} measure task success in code, web, and general settings. These measure
capability. We measure whether the tool implementations themselves are tested by the
projects that ship them---a software-engineering property of the codebase, orthogonal to
how well a model uses the tool.

Closest to our artefact is ToolEmu~\cite{toolemu}, which uses a language model to
emulate tool execution in order to surface agent risks without real side effects, and
AgentDojo~\cite{agentdojo}, which evaluates agents against prompt-injection attacks in
environments with stateful tools. Both establish that consequences of agent writes
deserve dedicated infrastructure. AgentCage differs in mechanism and in aim: our model
is deterministic code rather than an LM emulation, which trades breadth of coverage for
exact, inspectable state transitions and reproducible failures; and our target is the
developer's own test suite rather than a red-teaming evaluation.

We are not aware of prior work quantifying test coverage of write-capable tools across
open-source agent repositories, which is the gap Section~\ref{sec:parta} addresses.

\section{Part A: how much is tested?}
\label{sec:parta}

\subsection{Sampling frame}

We assembled \AcCandidates{} candidate repositories: recognisable Python agent
projects spanning frameworks, autonomous agents, coding agents, and browser agents.
The frame is a convenience sample, not a random sample of GitHub, and is restricted to
Python-primary repositories because every detector is a Python source pattern that
would silently under-detect on other languages.

This biases the result toward well-maintained projects, which if anything biases
measured coverage \emph{upward}: the long tail of small agent repositories is unlikely
to be better tested than CrewAI or LlamaIndex.

\subsection{Screening: what enters the denominator}

A repository enters the analysed set only if the detectors find at least one
\emph{write-capable tool definition}. Screening is mechanical: \AcScreenedOut{} of
\AcCandidates{} candidates were excluded by code, leaving \AcRepos{}. Screening out is
not a claim that a project performs no writes; it means the project exposes no
framework-idiomatic tool definition whose body performs a detectable mutation.
Section~\ref{sec:threats} treats this as the study's principal limitation.

\subsection{Detectors}

Tool definitions are recognised in three styles: \textbf{decorator}
(\texttt{@tool}, \texttt{@agent.tool}, \texttt{@mcp.tool()}, \texttt{@registry.action},
\texttt{@function\_tool}, \texttt{@kernel\_function}), \textbf{class} (a base named
\texttt{*BaseTool}, \texttt{*Toolkit}, \texttt{*ToolSpec}, \texttt{*BaseAction}), and
\textbf{schema} (a JSON tool schema pairing a \texttt{name} with \texttt{parameters} or
\texttt{input\_schema}). CLI decorators such as \texttt{@click.command} are deliberately
excluded, since matching them would classify ordinary command-line programs as agents.

A definition is write-capable if its body matches at least one non-ambiguous mutation
marker across ten categories: destructive HTTP verbs, filesystem mutation, file opens
in write mode, shell execution, SQL and ORM writes, cloud SDK writes, version-control
writes, outbound messaging, and payment operations. A bare HTTP \texttt{POST} is placed
in a separate tier and \emph{excluded}, because many read-only search and inference
endpoints are queried with \texttt{POST}.

Two rules prevent a definition from serving as its own evidence. First, markers are
matched only against the body \emph{after} the signature: without this,
\texttt{def delete\_file(path)} satisfies a pattern intended to catch \emph{calls} to
\texttt{delete\_file(}, and any tool whose name reads like a mutation is counted as one
on the strength of its own name. Second, every occurrence of the tool's own identifier
is neutralised before matching. Both rules were added after manual inspection revealed
false positives that these exact mechanisms had produced.

\subsection{Coverage criterion}

For each analysed repository we read its test sources and count a write tool as
\textbf{covered} if its identifier appears as a whole word anywhere in that corpus.

This bar is intentionally generous. It requires no assertion, no exercise of the
destructive path, and no execution: a single import or an incidental mention in an
unrelated test satisfies it. Coverage measured this way is therefore an \emph{upper
bound} on real coverage, and the uncovered fraction a \emph{lower bound}. A tool this
study calls untested is untested under the weakest available definition.

\paragraph{Unit of analysis.} We count \emph{tools}, not repositories. A repository-level
criterion---``has at least one test touching at least one write tool''---is nearly
vacuous: a project shipping 34 write tools with one of them referenced satisfies it.
Counting tools places \AcTools{} items in the denominator rather than \AcRepos{}.

\subsection{Results}

Table~\ref{tab:parta} reports per-repository counts for the run of \AcRunDate{}.
Across \AcRepos{} repositories, \AcUncovered{} of \AcTools{} write-capable tool
definitions (\AcUncoveredPct{}) are never referenced by any test.

% GENERATED by paper/make_tables.py -- do not edit by hand.
\begin{table}[t]
\centering
\small
\begin{tabular}{lrrrrl}
\toprule
Repository & Tools & Covered & Uncovered & Uncov. & Test scan \\
\midrule
agno-agi/agno & 34 & 23 & 11 & 32\% & capped \\
crewAIInc/crewAI & 32 & 19 & 13 & 41\% & full \\
camel-ai/camel & 30 & 29 & 1 & 3\% & full \\
run-llama/llama\_index & 13 & 13 & 0 & 0\% & capped \\
geekan/MetaGPT & 13 & 13 & 0 & 0\% & full \\
langchain-ai/langchain-community & 12 & 11 & 1 & 8\% & capped \\
TransformerOptimus/SuperAGI & 9 & 6 & 3 & 33\% & full \\
Upsonic/Upsonic & 6 & 4 & 2 & 33\% & capped \\
Significant-Gravitas/AutoGPT & 3 & 3 & 0 & 0\% & capped \\
griptape-ai/griptape & 3 & 3 & 0 & 0\% & capped \\
microsoft/autogen & 1 & 1 & 0 & 0\% & full \\
openai/openai-agents-python & 1 & 1 & 0 & 0\% & full \\
aiwaves-cn/agents & 1 & 0 & 1 & 100\% & full \\
\midrule
\textbf{Total} & \textbf{158} & \textbf{126} & \textbf{32} & \textbf{20.3\%} & \\
\bottomrule
\end{tabular}
\caption{Write-capable tool definitions and their test coverage, for the run of 2026-08-04. ``Test scan'' is \emph{full} where the repository's entire test suite was read and \emph{capped} where the per-repository read limit bound. Generated directly from \texttt{part\_a/results.json}.}
\label{tab:parta}
\end{table}


\paragraph{Sensitivity.} Restricting to the \AcSensRepos{} repositories whose entire
test suite was read---where a missed reference cannot be an artefact of sampling---gives
\AcSensUncovered{} of \AcSensTools{} tools uncovered, or \AcSensPct{}. The agreement
with the headline indicates the estimate is not driven by the read cap.

\paragraph{Repository-level view.} \AcRepoUntested{} of \AcRepos{} repositories
(\AcRepoUntestedPct{}) have no test touching any write tool: \AcUntestedRepoList{},
which ships no test files at all. We report this figure for completeness while noting
it rests on a single repository and is far less informative than the tool-level result.

\subsection{Detector revision history}
\label{sec:revisions}

The detectors were developed against the corpus, and each revision moved the headline.
We record them because the final number is not interpretable without them.

\begin{table}[t]
\centering
\small
\begin{tabular}{p{0.44\textwidth}p{0.44\textwidth}}
\toprule
Problem identified & Change and effect \\
\midrule
Path filter matched substrings, so \texttt{openai.py} matched an ``op'' alternative &
Match whole delimiter-separated path words; candidate files became genuine tool surfaces \\
File cap applied lexicographically, representing large repositories by whichever
directory sorted first & Rank by tool-likeness before applying the cap \\
Reading only files whose \emph{path} contained a tool keyword missed tools defined in
files such as \texttt{controller/service.py} & Make every non-test source file eligible;
\AcRepos{}$\rightarrow$17 repositories analysed \\
\texttt{@registry.action} and MCP-style decorators absent from the vocabulary & Added;
included in the above \\
\texttt{db\_write} matched bare \texttt{.save(}, catching a PIL image save & Removed bare
\texttt{.save(}; removed a false positive from the untested set \\
Patterns matched the definition line itself, classifying a tool as a write on the
strength of \texttt{def delete\_file(} & Strip signatures, neutralise self-reference;
17$\rightarrow$\AcRepos{} repositories \\
Repository-level binary coverage too weak, and after the above rested on one repository
& Switch unit of analysis to the tool; headline became \AcUncoveredPct{} of \AcTools{} \\
\bottomrule
\end{tabular}
\caption{Detector revisions. The final two rows removed findings that would have
\emph{raised} the reported figure.}
\label{tab:revisions}
\end{table}

\paragraph{Status of this analysis.} Part A is \textbf{exploratory}, not
pre-registered. The detectors were revised in response to inspecting their output, as
Table~\ref{tab:revisions} records. It should be read as a measurement with a documented
derivation rather than a confirmatory test of a hypothesis fixed in advance. The Part B
evaluation in Section~\ref{sec:partb} \emph{is} pre-registered and has not been run.

\subsection{Threats to validity}
\label{sec:threats}

\paragraph{Construct validity is the dominant limitation.} The study measures
\emph{framework-idiomatic tool definitions}. Agents with bespoke architectures---%
\texttt{aider}, \texttt{OpenHands}, \texttt{SWE-agent}, \texttt{gpt-engineer}---write to
disk and execute shell commands continuously but expose no \texttt{@tool} or
\texttt{BaseTool} surface, and therefore screen out. With \AcScreenedOut{} of
\AcCandidates{} candidates excluded, the denominator of \AcRepos{} reflects detector
coverage as much as it reflects the ecosystem. The correct reading of our headline is
``among tools we can mechanically identify,'' not ``among all agent write paths.''

\paragraph{Coverage inference.} Name reference in a test file is not execution. A
referenced tool may carry no assertion on its destructive path (inflating measured
coverage); an unreferenced tool may be exercised by an integration harness that never
names it (deflating it). We argue the first error dominates because the criterion is so
weak, which makes \AcUncoveredPct{} a floor.

\paragraph{Browser and GUI writes.} Agents that mutate state by clicking and submitting
forms are not detected, since no marker recognises a DOM interaction.

\paragraph{Frameworks versus applications.} Several screened-out projects ship the
\emph{mechanism} for defining tools rather than write tools themselves. Excluding them
is correct, but it skews the analysed set toward tool catalogues.

\paragraph{Temporal validity.} Files are read at \texttt{HEAD}. The measurement is of a
moving target; results carry a timestamp and continuous integration re-runs the study
weekly.

\section{Part B: AgentCage}
\label{sec:partb}

\subsection{Design}

\textbf{Interceptor.} Attaches through each HTTP library's own documented hook points
(\texttt{httpx} event hooks, \texttt{requests} response hooks) and records method, URL,
path, query, headers, body, and response status. It is strictly passive: it never
rewrites a request, injects a response, or short-circuits the transport, and a test
asserts that the caller receives exactly what the transport returned. Credential headers
are redacted at capture time so traces are committable artefacts.

\textbf{State model.} A deterministic in-memory model of the money-moving subset of the
Stripe API. Refunds increase \texttt{charge.amount\_refunded} and set
\texttt{charge.refunded} at parity with \texttt{amount\_captured}; a subsequent
\texttt{GET} observes the change. Refunds beyond the unrefunded amount are rejected with
the real error code and message; a fully refunded charge rejects further refunds;
idempotency keys replay the original response and reject reuse with altered parameters.
Identifiers are counter-derived and the clock is injectable, so runs are reproducible.

\subsection{Fidelity boundary}

Not modelled: payment-method validation, 3-D Secure, disputes, payouts, balance
transactions, webhooks, connected accounts, field expansion, pagination cursors, rate
limits, and Stripe's real identifier format. The model is sufficient to exercise
double-refund, over-refund, refund-after-refund, and read-after-write, and no more.
Stating this boundary matters: a mock's silent divergence from the real API is itself a
source of false confidence.

\subsection{Worked example}

The smallest complete demonstration, produced by \texttt{part\_b/demo\_double\_refund.py}
and committed as \texttt{traces/example\_double\_refund.json}:

\begin{lstlisting}
POST /v1/charges         -> 200   ch_00000001 created, 5000 usd
POST /v1/refunds         -> 200   amount_refunded 0 -> 5000
POST /v1/refunds         -> 400   charge_already_refunded    <- the retry
GET  /v1/charges/ch_...  -> 200   amount_refunded == 5000
\end{lstlisting}

Both \texttt{POST /v1/refunds} requests are byte-identical. A replay fixture answers the
second with 200 and the customer is refunded twice; the stateful model refuses it. The
captured trace carries \texttt{authorization: <redacted>}, which is what makes traces
publishable artefacts.

\subsection{Pre-registered evaluation}
\label{sec:prereg}

\textbf{Status: not run. No results exist for this section.} It is stated before the
experiment precisely so that it cannot change afterwards.

\paragraph{H1.} Agents that pass their existing replay-based suites exhibit write-path
defects when the same paths are exercised against a stateful model.

\paragraph{Sample and exclusions.} Drawn from the Part A analysed set, restricted to
repositories whose tools perform payment-like writes. Pre-specified exclusions: no
runnable suite; write tools requiring credentials we cannot substitute; tools that mutate
only local filesystem state, giving the interceptor no HTTP surface to observe. Any agent
added or dropped after its result is known will be reported as such.

\paragraph{Procedure.} For each agent: run its own suite unmodified (expected: passes;
an agent whose suite already fails is excluded and reported); re-run its write paths
through the interceptor against the stateful model; grade every request by the resolution
tiers; classify every divergence by the defect classes.

\paragraph{Resolution tiers.} Each captured request receives exactly one tier, with
precedence \textbf{T1 $>$ T2 $>$ T3 $>$ MISS}; the grader records which rule fired.

\begin{table}[h]
\centering
\small
\begin{tabular}{p{0.13\textwidth}p{0.78\textwidth}}
\toprule
Tier & Definition \\
\midrule
\textbf{T1} & \emph{Exact.} Same status code and same decision-relevant fields as the
real API, and the agent's control flow is identical. \\
\textbf{T2} & \emph{Behaviourally equivalent.} Status class and decision-relevant
semantics match; only incidental fields differ (identifier format, timestamps, unmodelled
fields). Control flow unchanged. \\
\textbf{T3} & \emph{Divergent but exercising.} The response differs from the real API's
in a way that still drives the agent down its write path---typically the model is
stricter. Recorded as a fidelity gap. \\
\textbf{MISS} & The model cannot answer (unmodelled endpoint, parameter, or state), or
answers such that the scenario cannot be evaluated. \\
\bottomrule
\end{tabular}
\caption{Resolution tiers. T1 and T2 require a documented expectation of the real API's
behaviour; absent such a reference a request is graded T3 or MISS, never T1.}
\label{tab:tiers}
\end{table}

\paragraph{Miss classification.} Every MISS is assigned exactly one class, and this is
the measurement that decides whether the approach is viable at all.
\textbf{MODELABLE}---the gap closes by writing more model (an unmodelled endpoint, field,
or error path), so the cost is engineering effort. \textbf{PRODUCTION-STATE-DEPENDENT}---%
the miss requires state that exists only in a real account: a specific pre-existing
entity, a live dispute, a provider-side asynchronous event, an inbound webhook. No amount
of modelling closes it without real credentials.

\paragraph{Defect classes.} Fixed in advance and not extendable after results are seen:
(1) \emph{duplicate write}---the same logical mutation applied twice; (2) \emph{unchecked
error}---a 4xx treated as success; (3) \emph{stale read}---a decision taken on state
cached from before a write; (4) \emph{over-refund}---a refund exceeding the unrefunded
amount. A genuine defect fitting none of these is recorded as \emph{unclassified},
reported separately, and does not count toward the primary outcome.

\paragraph{Outcomes.} \emph{Primary:} proportion of sampled agents exhibiting at least
one class-1--4 defect. \emph{Secondary:} distribution of resolution tiers across all
captured requests; MISS rate; the MODELABLE / PRODUCTION-STATE-DEPENDENT split; defects
per agent by class. Analysis is descriptive proportions with the denominator stated, and
no subgroup analysis beyond these breakdowns.

\paragraph{Falsification.} H1 is not supported if fewer than \textbf{20\%} of sampled
agents exhibit any class-1--4 defect. That is reported as a negative result, not
reframed.

\paragraph{Kill criterion.} If more than \textbf{50\%} of MISSes are
PRODUCTION-STATE-DEPENDENT, the stateful-mock approach cannot be made general by further
engineering, and the project pivots rather than continuing to build API models. Both
thresholds are fixed now, before any data exists, so that they cannot be moved later.

\section{Discussion}

The measurement and the artefact are related but independent. Even if AgentCage were
never adopted, \AcUncoveredPct{} of identifiable write-capable agent tools going
unreferenced by tests is a fact about how this software is currently built. And even a
perfect coverage figure would not establish that fixture-based tests \emph{detect}
write-path defects, which is why Part B is posed as a falsifiable claim rather than a
demonstration.

The honest summary of Part A is narrower than its headline: among tools we can
mechanically identify in a convenience sample of well-known Python agent projects, about
one in five is never named by a test. Whether that generalises to the long tail, to
non-Python ecosystems, or to bespoke agent architectures is unresolved, and the
screening rate of \AcScreenedOut{}/\AcCandidates{} is the clearest signal of how much
remains unmeasured.

\section{Availability}

Code, the sampling frame, complete per-repository results, and the detector revision
history are at \url{https://github.com/RithikSatarla/agentcage}. The study reproduces
with \texttt{python -m part\_a.run}, requires no credentials, and is re-run weekly in
continuous integration. Every statistic in this paper is generated from
\texttt{part\_a/results.json} rather than transcribed.

\appendix

\section{Pre-registration checklist}
\label{app:checklist}

Applies to the Part B evaluation of Section~\ref{sec:prereg}. Part A is exploratory and
makes no pre-registration claim.

\begin{table}[h]
\centering
\small
\begin{tabular}{p{0.62\textwidth}p{0.30\textwidth}}
\toprule
Item & Committed \\
\midrule
Hypothesis stated before data collection & Yes \\
Sampling frame and exclusions fixed in advance & Yes \\
Primary outcome specified, and single & Yes \\
Secondary outcomes enumerated & Yes \\
Grading rubric with explicit precedence rules & Yes (T1 $>$ T2 $>$ T3 $>$ MISS) \\
Defect classes closed to extension & Yes (four classes) \\
Falsification threshold numeric & Yes (20\%) \\
Kill criterion numeric & Yes (50\% of MISSes) \\
Analysis plan fixed & Yes \\
Raw evidence published & Yes (all traces committed) \\
Deviations to be logged with date and reason & Yes \\
\textbf{Results known at time of writing} & \textbf{No --- none exist} \\
\bottomrule
\end{tabular}
\caption{Pre-registration checklist for Part B.}
\end{table}

\section{Example captured trace}
\label{app:trace}

Abridged from the committed trace file, produced by running

\begin{center}
\texttt{python -m part\_b.demo\_double\_refund}
\end{center}

\noindent Four requests were captured, three of them writes. Credential headers are
replaced at capture time, which is what makes traces committable.

\begin{lstlisting}
{
  "count": 4,
  "traces": [
    { "method": "POST", "path": "/v1/charges",
      "headers": { "authorization": "<redacted>",
                   "content-type": "application/json" },
      "body": "{\"amount\": 5000, \"currency\": \"usd\"}",
      "status_code": 200 },

    { "method": "POST", "path": "/v1/refunds",
      "body": "{\"charge\": \"ch_00000001\", \"amount\": 5000}",
      "status_code": 200 },

    { "method": "POST", "path": "/v1/refunds",
      "body": "{\"charge\": \"ch_00000001\", \"amount\": 5000}",
      "status_code": 400,
      "response_body": "{\"error\": {\"type\": \"invalid_request_error\",
                         \"code\": \"charge_already_refunded\"}}" },

    { "method": "GET", "path": "/v1/charges/ch_00000001",
      "status_code": 200,
      "response_body": "{... \"amount_refunded\": 5000 ...}" }
  ]
}
\end{lstlisting}

The two refund requests are byte-identical in method, path, and body. Only a mock that
holds state can answer them differently, and answering them differently is what
correctness requires.

\begin{thebibliography}{99}

\bibitem{react} S. Yao, J. Zhao, D. Yu, N. Du, I. Shafran, K. Narasimhan, Y. Cao.
ReAct: Synergizing Reasoning and Acting in Language Models.
\emph{arXiv:2210.03629}, 2022.

\bibitem{toolformer} T. Schick, J. Dwivedi-Yu, R. Dessì, R. Raileanu, M. Lomeli,
L. Zettlemoyer, N. Cancedda, T. Scialom. Toolformer: Language Models Can Teach
Themselves to Use Tools. \emph{arXiv:2302.04761}, 2023.

\bibitem{gorilla} S. G. Patil, T. Zhang, X. Wang, J. E. Gonzalez. Gorilla: Large
Language Model Connected with Massive APIs. \emph{arXiv:2305.15334}, 2023.

\bibitem{toolllm} Y. Qin et al. ToolLLM: Facilitating Large Language Models to Master
16000+ Real-world APIs. \emph{arXiv:2307.16789}, 2023.

\bibitem{swebench} C. E. Jimenez, J. Yang, A. Wettig, S. Yao, K. Pei, O. Press,
K. Narasimhan. SWE-bench: Can Language Models Resolve Real-World GitHub Issues?
\emph{arXiv:2310.06770}, 2023.

\bibitem{webarena} S. Zhou et al. WebArena: A Realistic Web Environment for Building
Autonomous Agents. \emph{arXiv:2307.13854}, 2023.

\bibitem{agentbench} X. Liu et al. AgentBench: Evaluating LLMs as Agents.
\emph{arXiv:2308.03688}, 2023.

\bibitem{gaia} G. Mialon et al. GAIA: A Benchmark for General AI Assistants.
\emph{arXiv:2311.12983}, 2023.

\bibitem{toolemu} Y. Ruan et al. Identifying the Risks of LM Agents with an
LM-Emulated Sandbox. \emph{arXiv:2309.15817}, 2023.

\bibitem{agentdojo} E. Debenedetti et al. AgentDojo: A Dynamic Environment to Evaluate
Prompt Injection Attacks and Defenses for LLM Agents. \emph{arXiv:2406.13352}, 2024.

\end{thebibliography}

\end{document}
